% ---
% title: Settling a payment
% tags: continuous, density
% ---

\begin{problem}
	The time $X$, in hours, that a bank needs to settle an incoming
	international payment is a continuous random variable with density
	\[ f(x)=\begin{cases}
			c\,x(2-x), & 0\le x\le 2, \\
			0,         & \text{otherwise}.
		\end{cases} \]
	(a) Find the constant $c$.
	(b) Find $P(X>1)$.
	(c) Find $P(0.5<X<1.5)$.
	(d) What is $P(X=1)$?
\end{problem}

\begin{slideproblem}
	The time $X$ (hours) a bank needs to settle an incoming international
	payment has density $f(x)=c\,x(2-x)$ for $0\le x\le 2$, and $f(x)=0$
	otherwise.
	(a) Find $c$. (b) Find $P(X>1)$. (c) Find $P(0.5<X<1.5)$.
	(d) What is $P(X=1)$?
\end{slideproblem}

\begin{slides}
	\small
	\exstatement
	\pause
	\begin{columns}[T]
		\begin{column}{.62\textwidth}
			\sol Probability $=$ area under $f$; the total area is 1.\\[1mm]
			(a) $\displaystyle\int_0^2 c\,(2x-x^2)\,dx
				=c\Big[x^2-\frac{x^3}{3}\Big]_0^2=\frac43\,c=1$,
			\ so \ $\alert{c=\tfrac34}$.\\[1mm]
			\pause
			(b) $f$ is symmetric around $x=1$: \ $\alert{P(X>1)=\frac12}$.\\[1mm]
			\pause
			(c) $\displaystyle P(0.5<X<1.5)=\frac34\Big[x^2-\frac{x^3}{3}\Big]_{0.5}^{1.5}
				=\frac34\Big(\frac98-\frac5{24}\Big)=\alert{\frac{11}{16}}=0.6875$.\\[1mm]
			\pause
			(d) No area above a single point: \ $\alert{P(X=1)=0}$.
			So $<$ or $\le$ does not matter, and $f(1)=0.75$ is not a probability.
		\end{column}
		\begin{column}{.36\textwidth}
			\centering
			\begin{tikzpicture}[xscale=1.75,yscale=2.6]
				\draw[gray!40,very thin] (0,.25)--(2.2,.25) (0,.5)--(2.2,.5) (0,.75)--(2.2,.75);
				\node[left,font=\tiny] at (0,.25) {$0.25$};
				\node[left,font=\tiny] at (0,.5) {$0.5$};
				\node[left,font=\tiny] at (0,.75) {$0.75$};
				\fill[teal!22] (.5,0) -- plot[domain=.5:1.5,samples=40]
					(\x,{.75*\x*(2-\x)}) -- (1.5,0) -- cycle;
				\draw[teal] (.5,0)--(.5,{.75*.5*1.5}) (1.5,0)--(1.5,{.75*.5*1.5});
				\node[font=\tiny,text=teal!70!black] at (1,.3) {$0.6875$};
				\draw[->] (0,0) -- (2.35,0) node[right,font=\scriptsize] {$x$};
				\draw[->] (0,0) -- (0,.95) node[above,font=\scriptsize] {$f(x)$};
				\draw[navy,thick] plot[domain=0:2,samples=50] (\x,{.75*\x*(2-\x)});
				\draw[gold,thick,dashed] (1,0) -- (1,.75);
				\foreach \x in {0.5,1,1.5,2}
				\node[below,font=\tiny] at (\x,0) {$\x$};
			\end{tikzpicture}
		\end{column}
	\end{columns}
\end{slides}

\begin{solution}
	\idea A density $f$ has two properties: $f(x)\ge0$ for every $x$, and the
	total area under its graph is 1. Probabilities are areas under $f$:
	$P(a<X<b)=\int_a^b f(x)\,dx$. The total-area condition gives $c$; the other
	parts are areas under the graph.

	\step{(a)} Outside $[0,2]$ the density is $0$, so the total area is the
	integral over $[0,2]$ only. First integrate without the constant. Multiply
	out, $x(2-x)=2x-x^2$; an antiderivative is $x^2-\frac{x^3}{3}$. Evaluate it
	at the bounds:
	\[ \int_0^2(2x-x^2)\,dx=\Big[x^2-\frac{x^3}{3}\Big]_0^2
		=\Big(4-\frac83\Big)-(0-0)=\frac43. \]
	The total area must be 1, so
	\[ c\cdot\frac43=1,\qquad c=\frac34. \]
	Check the other property: for $0\le x\le2$ both $x$ and $2-x$ are
	nonnegative, so $f(x)=\frac34x(2-x)\ge0$. The density is
	$f(x)=\frac34x(2-x)$ on $[0,2]$: a parabola arch with its highest point
	$f(1)=\frac34$ at $x=1$.

	\step{(b)} The graph is symmetric around $x=1$: $f(1+t)=\frac34(1+t)(1-t)
	=f(1-t)$. So the area to the right of 1 equals the area to the left of 1.
	The two areas add up to 1, so each is $\frac12$:
	\[ P(X>1)=\frac12. \]
	By integration, with the same antiderivative:
	\[ P(X>1)=\int_1^2\frac34(2x-x^2)\,dx
		=\frac34\Big[x^2-\frac{x^3}{3}\Big]_1^2
		=\frac34\Big(\frac43-\frac23\Big)=\frac34\cdot\frac23=\frac12. \]

	\step{(c)} The probability is the shaded area in the picture:
	\[ P(0.5<X<1.5)=\int_{0.5}^{1.5}\frac34(2x-x^2)\,dx
		=\frac34\Big[x^2-\frac{x^3}{3}\Big]_{0.5}^{1.5}. \]
	Evaluate the antiderivative at the upper and at the lower bound:
	\[ 1.5^2-\frac{1.5^3}{3}=2.25-1.125=1.125=\frac98,\qquad
		0.5^2-\frac{0.5^3}{3}=0.25-\frac{0.125}{3}=\frac{5}{24}\approx0.2083. \]
	So
	\[ P(0.5<X<1.5)=\frac34\Big(\frac98-\frac5{24}\Big)
		=\frac34\cdot\frac{22}{24}=\frac{11}{16}=0.6875. \]
	The payment is settled between half an hour and one and a half hours after
	arrival with probability $0.6875$.
	\begin{center}
		\begin{tikzpicture}[xscale=2.4,yscale=3.6]
			\fill[teal!30] (.5,0) -- plot[domain=.5:1.5,samples=40]
				(\x,{.75*\x*(2-\x)}) -- (1.5,0) -- cycle;
			\node[teal!70!black] at (1,.3) {\small $\frac{11}{16}$};
			\draw[->] (-.2,0) -- (2.4,0) node[right] {\small $x$};
			\draw[->] (0,-.05) -- (0,.95) node[above] {\small $f(x)$};
			\draw[navy,thick] plot[domain=0:2,samples=50] (\x,{.75*\x*(2-\x)});
			\foreach \x in {0.5,1,1.5,2}
			\draw (\x,.02) -- (\x,-.02) node[below] {\scriptsize $\x$};
			\draw (.04,.75) -- (-.04,.75) node[left] {\scriptsize $0.75$};
		\end{tikzpicture}
	\end{center}

	\step{(d)} A probability is an area under $f$, and above a single point
	there is no area:
	\[ P(X=1)=\int_1^1 f(x)\,dx=0. \]
	This does not mean that settling in exactly one hour is impossible; it
	means that a single exact value carries no probability for a continuous
	random variable. As a consequence, $P(0.5<X<1.5)=P(0.5\le X\le1.5)$: for a
	continuous random variable it does not matter whether we write $<$ or
	$\le$.
	\begin{answer}
		(a) $c=\frac34$ \quad (b) $P(X>1)=\frac12$ \quad
		(c) $P(0.5<X<1.5)=\frac{11}{16}=0.6875$ \quad (d) $P(X=1)=0$
	\end{answer}
	\begin{remark}
		\textbf{Common mistake:} reading $f(1)=0.75$ as a probability. The value of
		a density is not a probability (a density can even be larger than 1);
		only areas under it are. Here $f(1)=0.75$ only says that values near
		1 hour are the most likely: $P(1<X<1+h)\approx f(1)\,h$ for a short
		interval of length $h$.
	\end{remark}
\end{solution}
